% !TEX program = lualatex
\documentclass{ltjsarticle}

\begin{document}
\title{測定板の位置}\author{}\date{2026年5月14日}\maketitle
\begin{abstract}板に刻んだ二本の線から、取り付け時の位置を読み取る。\end{abstract}
\tableofcontents
\section{目盛の読み方}\begin{quote}「下の線から順に読む。」\end{quote}
\begin{figure}[ht]\centering\fbox{\rule{0pt}{8mm}\rule{33mm}{0pt}}\caption{板の取り付け位置}\end{figure}
\section{高さ}
\begin{table}[ht]\centering\begin{tabular}{cr}目盛&高さ (mm)\\\hline 1&12\\2&18\end{tabular}\caption{目盛の高さ}\end{table}
\begin{thebibliography}{1}\bibitem{record}\emph{測定台帳}, 2025, 第2章.\end{thebibliography}
\end{document}
